\documentclass[a4paper]{article}

\usepackage{hyperref}
\usepackage[utf8]{inputenc}
\usepackage[a4paper, margin=3.5cm]{geometry}
\usepackage{graphicx}
\usepackage{url}
\usepackage[spanish,es-tabla]{babel}
\usepackage{fancyhdr}
\usepackage{float}
\usepackage{booktabs}
\usepackage{multirow}
\usepackage{minted}
\usepackage[nolist]{acronym}
\usepackage[
    type={CC},
    modifier={by-nc-sa},
    version={3.0},
]{doclicense}

\renewcommand{\headrulewidth}{0.6pt}
\renewcommand{\footrulewidth}{0.6pt}

\pagestyle{fancy}
\setlength\headheight{50pt}
\lhead{\includegraphics[height=1.5cm]{logos/upm_logo}}
\chead{Bases de Datos\\\vspace{.5em} Ejercicios procedimientos, funciones y triggers\\\vspace{-.1em}}
\rhead{\includegraphics[height=1.5cm]{logos/etsisi_logo}}
\lfoot{\textbf{Tema 5:} Gestión de bases de datos}
\cfoot{}
\rfoot{\thepage}

\parskip 1.1ex % paragraph spacing

\begin{acronym}
    \acro{FIA}{Federación Internacional de Automovilismo}
\end{acronym}

\begin{document}
\section*{Carreras de Formula 1}

Soluciones a los procedimientos, funciones y triggers:

\begin{enumerate}
    \item 
    \begin{minted}[fontsize=\scriptsize, frame=lines, xleftmargin=\parindent, breaklines, framesep=2mm]{sql}
DELIMITER $$
CREATE PROCEDURE `getRacesInAYear`(IN `year` INT)
BEGIN
	SELECT races.name, COUNT(DISTINCT results.constructorId) AS numConstructors
    FROM results
    	INNER JOIN races ON races.raceId = results.raceId
    WHERE races.year = year
    GROUP BY races.raceId;
END$$
DELIMITER ;
    \end{minted}
    
    \item
    \begin{minted}[fontsize=\scriptsize, frame=lines, xleftmargin=\parindent, breaklines, framesep=2mm]{SQL}
DELIMITER $$
CREATE PROCEDURE getsOnRaceMessages(IN cod VARCHAR(3), OUT msg VARCHAR(200))
BEGIN
	CASE cod
    WHEN 'E01' THEN SET msg = 'Error en la presión de las ruedas';
    WHEN 'E02' THEN SET msg = 'Pinchazo';
    WHEN 'E03' THEN SET msg = 'Temperatura alta en el motor';
    WHEN 'E04' THEN SET msg = 'Frenos sobre-calentados';
    WHEN 'E05' THEN SET msg = 'Error presión del aceite';
    ELSE SET msg = 'Error de comando';
    END CASE;
END$$
DELIMITER ;
    \end{minted}
    
    \item
    \begin{minted}[fontsize=\scriptsize, frame=lines, xleftmargin=\parindent, breaklines, framesep=2mm]{SQL}
DELIMITER $$
CREATE PROCEDURE getDriversByNationality (IN nat VARCHAR(250))
BEGIN
    SELECT *
    FROM drivers
    WHERE nationality = nat;
END$$
DELIMITER ;
    \end{minted}
    
    \item
    \begin{minted}[fontsize=\scriptsize, frame=lines, xleftmargin=\parindent, breaklines, framesep=2mm]{SQL}
DELIMITER $$
CREATE FUNCTION puntosCampeon (year INTEGER)
RETURNS DECIMAL(10,2)
DETERMINISTIC
BEGIN
	DECLARE points DECIMAL(10,2);

    SELECT MAX(T.totalPoints) INTO points
    FROM (SELECT SUM(results.points) AS totalPoints
          FROM results
              INNER JOIN races ON races.raceId = results.raceId
          WHERE races.year = year
          GROUP BY results.driverId) AS T;
          
    RETURN (points);
END$$
DELIMITER ;
    \end{minted}
    
    \item
    \begin{minted}[fontsize=\scriptsize, frame=lines, xleftmargin=\parindent, breaklines, framesep=2mm]{SQL}
DELIMITER $$
CREATE FUNCTION mediaPuntosConstructor (constructor VARCHAR(200))
RETURNS DECIMAL(10,2)
DETERMINISTIC
BEGIN
	DECLARE points DECIMAL(10,2);

    SELECT AVG(T.totalPoints) INTO points
    FROM (SELECT SUM(results.points) AS totalPoints
          FROM results
              INNER JOIN races ON races.raceId = results.raceId
              INNER JOIN constructors ON constructors.constructorId = results.constructorId
          WHERE constructors.name = constructor
          GROUP BY races.year) AS T;
          
    RETURN (points);
END$$
DELIMITER ;
    \end{minted}
    
    \item
    \begin{minted}[fontsize=\scriptsize, frame=lines, xleftmargin=\parindent, breaklines, framesep=2mm]{SQL}
DELIMITER $$
CREATE FUNCTION añosEnActivo (id INTEGER)
RETURNS INTEGER
DETERMINISTIC
BEGIN
	DECLARE años INTEGER;

    SELECT COUNT(DISTINCT races.year) INTO años 
    FROM results
    	INNER JOIN races ON races.raceId = results.raceId
    WHERE results.driverId = id;
          
    RETURN (años);
END$$
DELIMITER ;
    \end{minted}
    
    \item
    \begin{minted}[fontsize=\scriptsize, frame=lines, xleftmargin=\parindent, breaklines, framesep=2mm]{SQL}
ALTER TABLE drivers ADD COLUMN añosEnActivo INTEGER NULL;

DELIMITER $$
CREATE PROCEDURE actualizarAñosEnActivo ()
BEGIN
    DECLARE done INTEGER DEFAULT FALSE;
    DECLARE id INTEGER;
    
    DECLARE cur CURSOR FOR SELECT driverId FROM drivers;
    DECLARE CONTINUE HANDLER FOR NOT FOUND SET done = TRUE;
    OPEN cur;
    read_loop: LOOP
        FETCH cur INTO id;
        
        UPDATE drivers
            SET drivers.añosEnActivo = añosEnActivo(id)
            WHERE driverId = id;
            
        IF done THEN
            LEAVE read_loop;
        END IF;
    END LOOP;
    CLOSE cur;
END$$
DELIMITER ;
    \end{minted} 
    
    \item 
    \begin{minted}[fontsize=\scriptsize, frame=lines, xleftmargin=\parindent, breaklines, framesep=2mm]{SQL}    
-------------TABLA results ----------
DELIMITER $$
CREATE TRIGGER noMasDeDosPilotos
BEFORE INSERT ON results
FOR EACH ROW
BEGIN
    DECLARE numResults INTEGER;
    DECLARE numDrivers INTEGER;
    
    SELECT COUNT(*) INTO numResults
    FROM results
    WHERE constructorId = NEW.constructorId
      AND driverId = NEW.driverId;
      
    IF numResults = 0 THEN
        SELECT COUNT(DISTINCT driverId) INTO numDrivers
        FROM results
        WHERE constructorId = NEW.constructorId;

        IF numDrivers > 1 THEN
            SIGNAL SQLSTATE '03000'
            SET MESSAGE_TEXT = 'Error: no puede haber equipos con más de dos pilotos';
        END IF;
    END IF;
END$$
DELIMITER ;
-------------TABLA qualifying ----------
DELIMITER $$
CREATE TRIGGER noMasDeDosPilotos
BEFORE INSERT ON qualifying
FOR EACH ROW
BEGIN
    DECLARE numResults INTEGER;
    DECLARE numDrivers INTEGER;
    
    SELECT COUNT(*) INTO numResults
    FROM qualifying
    WHERE constructorId = NEW.constructorId
      AND driverId = NEW.driverId;
      
    IF numResults = 0 THEN
        SELECT COUNT(DISTINCT driverId) INTO numDrivers
        FROM qualifying
        WHERE constructorId = NEW.constructorId;

        IF numDrivers > 1 THEN
            SIGNAL SQLSTATE '03000'
            SET MESSAGE_TEXT = 'Error: no puede haber equipos con más de dos pilotos';
        END IF;
    END IF;
END$$
DELIMITER ;
    \end{minted}
    
    \item 
    \begin{minted}[fontsize=\scriptsize, frame=lines, xleftmargin=\parindent, breaklines, framesep=2mm]{SQL}
CREATE TABLE crashes (
   	crashId INTEGER UNIQUE NOT NULL AUTO_INCREMENT,
    driverId INT NOT NULL,
    description VARCHAR(250) DEFAULT NULL,
  	PRIMARY KEY (crashId),
	CONSTRAINT
		FOREIGN KEY (driverId)
		REFERENCES drivers (driverId)
)

DELIMITER $$
CREATE TRIGGER resgistrarAccidentes
AFTER INSERT ON results
FOR EACH ROW
BEGIN
    IF NEW.statusId = 3 OR NEW.statusId = 4 THEN
    	INSERT INTO crashes (driverId, description) VALUES (NEW.driverId, 'blah blah blah');
    END IF;
END$$
DELIMITER ;            
            
    \end{minted}
    
    \item 
    \begin{minted}[fontsize=\scriptsize, frame=lines, xleftmargin=\parindent, breaklines, framesep=2mm]{SQL}
-- Para la solución se asume que "estar" en un equipo es participar en una carrera (tabla results)

DELIMITER $$
CREATE TRIGGER no_more_than_two_teams_in_a_year
BEFORE INSERT ON results
FOR EACH ROW
BEGIN
    DECLARE y YEAR;
    DELCARE num_teams INTEGER;

    SELECT year INTO y FROM races WHERE raceId = NEW.raceId;

    SELECT COUNT(DISTINCT constructorId) INTO num_teams
    FROM results
        INNER JOIN races ON races.raceId = results.raceId
    WHERE races.year = y
      AND driverId = NEW.driverId;

    IF num_teams >= 2 THEN 
        SIGNAL SQLSTATE '03000'
        SET MESSAGE_TEXT = 'Error: no se puede competir con más de dos equipos en un año';
    END IF;
END$$
DELIMITER ;

DELIMITER $$
CREATE TRIGGER no_more_than_two_teams_in_a_year
BEFORE UPDATE ON results
FOR EACH ROW
BEGIN
    DECLARE y YEAR;
    DELCARE num_teams INTEGER;

    IF NEW.driverId <> OLD.driverId OR NEW.constructorId <> OLD.constructorId THEN
        SELECT year INTO y FROM races WHERE raceId = NEW.raceId;

        SELECT COUNT(DISTINCT constructorId) INTO num_teams
        FROM results
            INNER JOIN races ON races.raceId = results.raceId
        WHERE races.year = y
        AND driverId = NEW.driverId;

        IF num_teams >= 2 THEN 
            SIGNAL SQLSTATE '03000'
            SET MESSAGE_TEXT = 'Error: no se puede competir con más de dos equipos en un año';
        END IF;
    END IF;
END$$
DELIMITER ;
    \end{minted}

    \item 
    \begin{minted}[fontsize=\scriptsize, frame=lines, xleftmargin=\parindent, breaklines, framesep=2mm]{SQL}
DROP PROCEDURE pilots_win_home_constructor_year;

DELIMITER $$
CREATE PROCEDURE pilots_win_home_constructor_year(IN year_win INTEGER)
BEGIN
    SELECT DISTINCT drivers.forename, drivers.surname, circuits.name
    FROM drivers INNER JOIN results ON drivers.driverId = results.driverId INNER JOIN races ON results.raceId=races.raceId INNER JOIN constructors ON results.constructorId=constructors.constructorId INNER JOIN circuits ON circuits.circuitId=races.circuitId 
    WHERE results.positionOrder=1 AND races.year=year_win AND drivers.nationality=constructors.nationality;
END$$
DELIMITER ;
    \end{minted}
    
    \item 
    \begin{minted}[fontsize=\scriptsize, frame=lines, xleftmargin=\parindent, breaklines, framesep=2mm]{SQL}
DELIMITER //
CREATE PROCEDURE allPodiumPositions(IN anyo INTEGER)
BEGIN
SELECT forename, surname
FROM drivers D
WHERE driverId IN(SELECT driverId
                  FROM results JOIN races ON results.raceId=races.raceId
                  WHERE positionOrder=1
                  AND year=anyo)
AND driverId IN(SELECT driverId
                  FROM results JOIN races ON results.raceId=races.raceId
                  WHERE positionOrder=2
                  AND year=anyo)
AND driverId IN(SELECT driverId
                  FROM results JOIN races ON results.raceId=races.raceId
                  WHERE positionOrder=3
                  AND year=anyo);
END// 
DELIMITER ;    
    \end{minted}
    
    \item 
    \begin{minted}[fontsize=\scriptsize, frame=lines, xleftmargin=\parindent, breaklines, framesep=2mm]{SQL}
DELIMITER $$
CREATE FUNCTION diffPoints (driver1 INTEGER, driver2 INTEGER)
RETURNS DOUBLE
DETERMINISTIC
BEGIN
    DECLARE points1 DOUBLE;
    DECLARE points2 DOUBLE;

    SELECT SUM(points) INTO points1
    FROM results
    WHERE driverId = driver1;
    
    SELECT SUM(points) INTO points2
    FROM results
    WHERE driverId = driver2;
          
    RETURN (points1 - points2);
END$$
DELIMITER ;    
    \end{minted}
    
    \item 
    \begin{minted}[fontsize=\scriptsize, frame=lines, xleftmargin=\parindent, breaklines, framesep=2mm]{SQL}
CREATE TABLE sponsors(
    sponsorId INTEGER UNIQUE NOT NULL,
    name VARCHAR(50) NOT NULL,
    type VARCHAR(20),
    amount INTEGER NOT NULL,
    raceId INTEGER,
    PRIMARY KEY(sponsorId),
    CONSTRAINT
    FOREIGN KEY(raceId)
    REFERENCES formula1.races(raceId))

DELIMITER //
CREATE TRIGGER check_spon1 BEFORE INSERT ON sponsors
FOR EACH ROW
BEGIN
IF NEW.amount>5000000 THEN SET NEW.type='Oficial';
ELSE SET NEW.type='Co-oficial';
END IF;
END//

DELIMITER //
CREATE TRIGGER check_spon2 BEFORE UPDATE ON sponsors
FOR EACH ROW
BEGIN
IF NEW.amount>5000000 THEN SET NEW.type='Oficial';
ELSE SET NEW.type='Co-oficial';
END IF;
END//
DELIMITER ;    
    \end{minted}
    
    \item 
    \begin{minted}[fontsize=\scriptsize, frame=lines, xleftmargin=\parindent, breaklines, framesep=2mm]{SQL}
DELIMITER //
CREATE PROCEDURE `getsConstructoresYPilotos` (año year)
BEGIN
    SELECT constructors.name, drivers.surname, SUM(results.points)
    FROM constructors JOIN results ON constructors.constructorId = results.constructorID
		JOIN races ON results.raceId = races.raceId
        JOIN drivers ON results.driverId = drivers.driverId
	WHERE races.`year`= año
    GROUP BY constructors.name, drivers.surname;
END//  
DELIMITER ;
    \end{minted}
    
    \item 
    \begin{minted}[fontsize=\scriptsize, frame=lines, xleftmargin=\parindent, breaklines, framesep=2mm]{SQL}
DELIMITER $$
CREATE PROCEDURE getNumberOfVictories (IN type VARCHAR(20)) 
BEGIN
    IF type = 'nationality' THEN
        SELECT drivers.nationality, COUNT(*) AS numVictorires
        FROM results
            INNER JOIN drivers ON drivers.driverId = results.driverId
        WHERE results.positionOrder = 1
        GROUP BY drivers.nationality
        ORDER BY numVictorires DESC;
    ELSE 
        SELECT constructors.name, COUNT(*) AS numVictorires
        FROM results
            INNER JOIN constructors ON constructors.constructorId = results.constructorId
        WHERE results.positionOrder = 1
        GROUP BY constructors.name
        ORDER BY numVictorires DESC;
    END IF;
END$$
DELIMITER ;   
    \end{minted}
\end{enumerate}

\vspace{2em}
\hrule
\doclicenseThis

\end{document}
